%%%%%%%%%%%%%%%%%%%%%%%%%%%%%%%%%%%%%%%%%
% FRI Data Science_report LaTeX Template
% Version 1.0 (28/1/2020)
% 
% Jure Demšar (jure.demsar@fri.uni-lj.si)
%
% Based on MicromouseSymp article template by:
% Mathias Legrand (legrand.mathias@gmail.com) 
% With extensive modifications by:
% Antonio Valente (antonio.luis.valente@gmail.com)
%
% License:
% CC BY-NC-SA 3.0 (http://creativecommons.org/licenses/by-nc-sa/3.0/)
%
%%%%%%%%%%%%%%%%%%%%%%%%%%%%%%%%%%%%%%%%%


%----------------------------------------------------------------------------------------
%	PACKAGES AND OTHER DOCUMENT CONFIGURATIONS
%----------------------------------------------------------------------------------------
\documentclass[fleqn,moreauthors,10pt]{ds_report}
\usepackage[english]{babel}

\graphicspath{{fig/}}




%----------------------------------------------------------------------------------------
%	ARTICLE INFORMATION
%----------------------------------------------------------------------------------------

% Header
\JournalInfo{FRI Natural language processing course 2025}

% Interim or final report
\Archive{Project report} 
%\Archive{Final report} 

% Article title
\PaperTitle{AI chat assistant for law domain}

% Authors (student competitors) and their info
\Authors{Tibor Koderman, Jerneja Krajcar, Blaž Mikec}

% Advisors
\affiliation{\textit{Advisors: Slavko Žitnik}}

% Keywords
\Keywords{Law domain, AI chatbot, LLMD4LH, }
\newcommand{\keywordname}{Keywords}


%----------------------------------------------------------------------------------------
%	ABSTRACT
%----------------------------------------------------------------------------------------

\Abstract{
The abstract goes here. The abstract goes here. The abstract goes here. The abstract goes here. The abstract goes here. The abstract goes here. The abstract goes here. The abstract goes here. The abstract goes here. The abstract goes here. The abstract goes here. The abstract goes here. The abstract goes here. The abstract goes here. The abstract goes here. The abstract goes here. The abstract goes here. The abstract goes here. The abstract goes here. The abstract goes here. The abstract goes here. The abstract goes here. The abstract goes here. The abstract goes here. The abstract goes here. The abstract goes here.
}

%----------------------------------------------------------------------------------------

\begin{document}

% Makes all text pages the same height
\flushbottom 

% Print the title and abstract box
\maketitle 

% Removes page numbering from the first page
\thispagestyle{empty} 

%----------------------------------------------------------------------------------------
%	ARTICLE CONTENTS
%----------------------------------------------------------------------------------------

\section*{Introduction}
	These latex files are intended to serve as a the template for the NLP course at FRI.  The template is adapted from the FRI Data Science Project Competition. template  If you find mistakes in the template or have problems using it, please consult Jure Demšar (\href{mailto:jure.demsar@fri.uni-lj.si}{jure.demsar@fri.uni-lj.si}).
	
	In the Introduction section you should write about the relevance of your work (what is the purpose of the project, what will we solve) and about related work (what solutions for the problem already exist). Where appropriate, reference scientific work conducted by other researchers. For example, the work done by Demšar et al. \cite{Demsar2016BalancedMixture} is very important for our project. The abbreviation et al. is for et alia, which in latin means and others, we use this abbreviation when there are more than two authors of the work we are citing. If there are two authors (or if there is a single author) we just write down their surnames. For example, the work done by Demšar and Lebar Bajec \cite{Demsar2017LinguisticEvolution} is also important for successful completion of our project.


%------------------------------------------------

\section*{Existing Work}

In recent years, several solutions have emerged that aim to simplify access to legal information
and provide assistance in understanding national legislation. In the context of Slovenian law,
we identified a small number of relevant platforms that partially address this problem.

One notable example is \textit{Pravko.si}, which functions as a conversational assistant
for legal questions. \cite{Pravko.si} The system provides relatively precise answers and, importantly,
supports its responses with references to relevant legal acts and related materials.
This increases transparency and allows users to verify the correctness of the information.
Additionally, the assistant is capable of simplifying legal language to some extent,
making it more accessible to non-expert users. However, we observed that the system struggles
in cases where the requested information is not available in its knowledge base,
leading to unresponsive behavior instead of providing a fallback explanation or clarification.

Another platform is \textit{MojMaliPravnik.net}, which operates more as a question-and-answer portal
than an autonomous AI assistant. \cite{MojMaliPravnik.net} Users can submit legal questions, and answers are either provided
automatically or by human experts. Free questions and answers are published publicly (in anonymized form),
creating a shared knowledge base that can benefit other users with similar issues.
The automatically generated responses often resemble excerpts from legal documents,
which may limit their readability and practical usefulness. For more complex inquiries,
the platform offers paid individual legal advice, which is handled by professionals rather
than automated systems.

From this overview, we observe that existing solutions either focus on conversational interaction
with limited robustness, or rely heavily on human expertise with minimal automation.
Furthermore, there is a lack of systems that combine accurate referencing of formal legal acts
with reliable handling of missing information and consistently understandable explanations.

We also observe that general-purpose language models are ill-suited for Slovenian legal question answering.
Trained on broad multilingual data, they lack specialized knowledge of Slovenian statutes and often
hallucinate plausible but false information, such as inventing statute numbers.
In practice, models like ChatGPT have been found to generate completely fabricated case citations
and legal quotes.

This gap motivates the development of our proposed AI assistant, which aims to provide grounded,
explainable, and robust answers based directly on formal Slovenian legal sources while maintaining
accessibility for general users.

\section*{Corpus Analysis}

For the development of our system, we selected \textit{COLESLAW 1.0}, a recent collection of
Slovenian legal texts published in the CLARIN.SI repository. \cite{COLESLAW1.0} The corpus contains
547,799 texts and 771,931,725 words distributed across 12 JSONL files. At the highest level, it
combines four complementary legal domains: the PISRS legal information system, the
\textit{SodnaPraksa} collection of court decisions, the Constitutional Court collection
(\textit{USRS}), and texts from the \textit{Uradni list}. This is useful for our task because the
dataset does not focus only on laws or only on court decisions, but instead combines several types
of legal texts in one place.

A closer look at the internal structure shows that the corpus is not balanced across domains, which
is important for downstream modeling choices. The largest share by number of documents comes from
\textit{Uradni list}, which contributes 274,531 texts, or just over half of the corpus, while
\textit{SodnaPraksa} contributes 217,409 texts. In terms of word count, however,
\textit{SodnaPraksa} becomes the dominant source with 352.3 million words (45.6\% of all words),
followed by \textit{Uradni list} with 296.7 million words (38.4\%). PISRS is much smaller in the
number of documents, with 40,749 texts, but still contributes 96.8 million words, which indicates
that its documents are on average longer and denser than many gazette entries. The USRS collection
adds a further 15,110 constitutional court documents and 26.2 million words.

The corpus also contains noticeable differences between subcollections. For example,
\textit{sp\_courts} is the largest individual part of the dataset, while \textit{ul-razglasni}
contains shorter announcement-style texts. Some legislative collections in PISRS are much longer,
especially the register of regulations and draft regulations. Because of this, the corpus includes
both short notices and long legal documents. In our opinion, this makes COLESLAW 1.0 a good
starting point for building a Slovenian legal assistant, since it is large, diverse, and based on
official public sources.

%\subsection*{Equations}

%You can write equations inline, e.g. $\cos\pi=-1$, $E = m \cdot c^2$ and $\alpha$, or you can include them as separate objects. The Bayes’s rule is stated mathematically as:

%\begin{equation}
%	P(A|B) = \frac{P(B|A)P(A)}{P(B)},
%	\label{eq:bayes}
%\end{equation}

%where $A$ and $B$ are some events. You can also reference it -- the equation \ref{eq:bayes} describes the Bayes's rule.

%\subsection*{Lists}

%We can insert numbered and bullet lists:

% the [noitemsep] option makes the list more compact
%\begin{enumerate}[noitemsep]
%	\item First item in the list.
%	\item Second item in the list.
%	\item Third item in the list.
%\end{enumerate}

%\begin{itemize}[noitemsep]
%	\item First item in the list.
%	\item Second item in the list.
%	\item Third item in the list.
%\end{itemize}

%We can use the description environment to define or describe key terms and phrases.

%\begin{description}
%	\item[Word] What is a word?.
%	\item[Concept] What is a concept?
%	\item[Idea] What is an idea?
%\end{description}

%\subsection*{Figures}

%You can insert figures that span over the whole page, or over just a single column. The first one, \figurename~\ref{fig:column}, is an example of a figure that spans only across one of the two columns in the report.

%\begin{figure}[ht]\centering
%	\includegraphics[width=\linewidth]{single_column.pdf}
%	\caption{\textbf{A random visualization.} This is an example of a figure that spans only across one of the two columns.}
%	\label{fig:column}
%\end{figure}

%On the other hand, \figurename~\ref{fig:whole} is an example of a figure that spans across the whole page (across both columns) of the report.

% \begin{figure*} makes the figure take up the entire width of the page
%\begin{figure*}[ht]\centering
%	\includegraphics[width=\linewidth]{whole_page.pdf}
%	\caption{\textbf{Visualization of a Bayesian hierarchical model.} This is an example of a figure that spans the whole width of the report.}
%	\label{fig:whole}
%\end{figure*}


%\subsection*{Tables}

%Use the table environment to insert tables.

%\begin{table}[hbt]
%	\caption{Table of grades.}
%	\centering
%	\begin{tabular}{l l | r}
%		\toprule
%		\multicolumn{2}{c}{Name} \\
%		\cmidrule(r){1-2}
%		First name & Last Name & Grade \\
%		\midrule
%		John & Doe & $7.5$ \\
%		Jane & Doe & $10$ \\
%		Mike & Smith & $8$ \\
%		\bottomrule
%	\end{tabular}
%	\label{tab:label}
%\end{table}


%\subsection*{Code examples}

%You can also insert short code examples. You can specify them manually, or insert a whole file with code. Please avoid inserting long code snippets, advisors will have access to your repositories and can take a look at your code there. If necessary, you can use this technique to insert code (or pseudo code) of short algorithms that are crucial for the understanding of the manuscript.

%\lstset{language=Python}
%\lstset{caption={Insert code directly from a file.}}
%\lstset{label={lst:code_file}}
%\lstinputlisting[language=Python]{code/example.py}

%\lstset{language=R}
%\lstset{caption={Write the code you want to insert.}}
%\lstset{label={lst:code_direct}}
%\begin{lstlisting}
%import(dplyr)
%import(ggplot)

%ggplot(diamonds,
%	   aes(x=carat, y=price, color=cut)) +
%  geom_point() +
%  geom_smooth()
%\end{lstlisting}

%----------------------------------------------------------------------------------------
%	REFERENCE LIST
%----------------------------------------------------------------------------------------
\bibliographystyle{unsrt}
\bibliography{report}


\end{document}
